
\documentclass{article}
\usepackage{colortbl}
\usepackage{makecell}
\usepackage{multirow}
\usepackage{supertabular}

\begin{document}

\newcounter{utterance}

\centering \large Interaction Transcript for game `hot\_air\_balloon', experiment `air\_balloon\_survival\_en\_complexity\_easy', episode 1 with deepseek{-}chat{-}v3.1{-}t1.0.
\vspace{24pt}

{ \footnotesize  \setcounter{utterance}{1}
\setlength{\tabcolsep}{0pt}
\begin{supertabular}{c@{$\;$}|p{.15\linewidth}@{}p{.15\linewidth}p{.15\linewidth}p{.15\linewidth}p{.15\linewidth}p{.15\linewidth}}
    \# & $\;$A & \multicolumn{4}{c}{Game Master} & $\;\:$B\\
    \hline

    \theutterance \stepcounter{utterance}  
    & & \multicolumn{4}{p{0.6\linewidth}}{
        \cellcolor[rgb]{0.9,0.9,0.9}{
            \makecell[{{p{\linewidth}}}]{
                \texttt{\tiny{[P1$\langle$GM]}}
                \texttt{You are participating in a collaborative negotiation game.} \\
\\ 
\texttt{Together with another participant, you must agree on a single set of items that will be kept. Each of you has your own view of how much each item matters to you (importance). You do not know how the other participant values the items. Additionally, you are given the effort each item demands.} \\
\texttt{You may only agree on a set if the total effort of the selected items does not exceed a shared limit:} \\
\\ 
\texttt{LIMIT = 8170} \\
\\ 
\texttt{Here are the individual item effort values:} \\
\\ 
\texttt{Item effort = \{"B21": 594, "B00": 551, "A08": 634, "A27": 204, "B94": 557, "B53": 440, "A23": 734, "C65": 893, "A89": 246, "B67": 592, "B45": 143, "C34": 568, "C12": 472, "A05": 401, "B60": 729, "C92": 201, "B69": 85, "A84": 641, "C06": 958, "B12": 80, "A16": 157, "B85": 806, "C35": 684, "B10": 59, "C15": 31, "B09": 765, "B16": 416, "C81": 392, "A58": 428, "C24": 699, "A30": 141, "A66": 606, "B75": 612, "B46": 133, "C70": 689\}} \\
\\ 
\texttt{Here is your personal view on the importance of each item:} \\
\\ 
\texttt{Item importance values = \{"B21": 138, "B00": 583, "A08": 868, "A27": 822, "B94": 783, "B53": 65, "A23": 262, "C65": 121, "A89": 508, "B67": 780, "B45": 461, "C34": 484, "C12": 668, "A05": 389, "B60": 808, "C92": 215, "B69": 97, "A84": 500, "C06": 30, "B12": 915, "A16": 856, "B85": 400, "C35": 444, "B10": 623, "C15": 781, "B09": 786, "B16": 3, "C81": 713, "A58": 457, "C24": 273, "A30": 739, "A66": 822, "B75": 235, "B46": 606, "C70": 105\}} \\
\\ 
\texttt{Goal:} \\
\\ 
\texttt{Your goal is to negotiate a shared set of items that benefits you as much as possible (i.e., maximizes total importance to YOU), while staying within the LIMIT. You are not required to make a PROPOSAL in every message {-} you can simply negotiate as well. All tactics are allowed!} \\
\\ 
\texttt{Interaction Protocol:} \\
\\ 
\texttt{You may only use the following structured formats in a message:} \\
\\ 
\texttt{PROPOSAL: \{'A', 'B', 'C', …\}} \\
\texttt{Propose keeping exactly those items.} \\
\texttt{REFUSE: \{'A', 'B', 'C', …\}} \\
\texttt{Explicitly reject opponent's proposal.} \\
\texttt{ARGUMENT: \{'...'\}} \\
\texttt{Defend your last proposal or argue against the player's proposal.} \\
\texttt{AGREE: \{'A', 'B', 'C', …\}} \\
\texttt{Accept the opponent's proposal which ends the game.} \\
\texttt{STRATEGIC REASONING: \{'...'\}} \\
\texttt{	Describe your strategic reasoning or anticipation explaining your choice of action. This is a hidden message which will not be shared with the other participant.} \\
\\ 
\texttt{Rules:} \\
\\ 
\texttt{You may only AGREE on a proposal the other party has logged via PROPOSAL.} \\
\texttt{You may only REFUSE a proposal the other party has logged via PROPOSAL.} \\
\texttt{Total effort of any PROPOSAL or AGREE set must be ≤ LIMIT.} \\
\texttt{Do NOT reveal your hidden importance scores.} \\
\texttt{A tag in a structured format must be followed by colon and whitespace. The argument must be a python set containing 0 or more strings.} \\
\texttt{So, it must be of the form TAG: \{...\}} \\
\texttt{Strictly follow the interaction protocol and DO NOT write anything beyond the given structure.} \\
\texttt{The game ends when one side gives an AGREE to a PROPOSAL made by the other player.} \\
\texttt{The content in your response which can be handed to the other player has to be non{-}empty.} \\
\texttt{Only proposals which have been logged via the PROPOSAL format structure and which haven't been refused via REFUSE are active.} \\
\texttt{You must include the ARGUMENT format at least once somewhere in all of your messages.} \\
\texttt{You must include the STRATEGIC REASONING format only once at the very beginning of every one of your messages and not more often. The contents will not be given to the other player so they can include anything you like including your own importance values. Here you should reason step by step to come up with you next move.} \\
            }
        }
    }
    & \\ \\

    \theutterance \stepcounter{utterance}  
    & \multicolumn{4}{p{0.6\linewidth}}{
        \cellcolor[rgb]{0.9,0.9,0.9}{
            \makecell[{{p{\linewidth}}}]{
                \texttt{\tiny{[P1$\rangle$GM]}}
                \texttt{STRATEGIC REASONING: \{'I need to maximize my total importance, which is 138+583+868+822+783+65+262+121+508+780+461+484+668+389+808+215+97+500+30+915+856+400+444+623+781+786+3+713+457+273+739+822+235+606+105 = 19879. However, we have an effort limit of 8170. I should calculate the effort{-}to{-}importance ratio for each item to prioritize items that give high importance per unit effort. Let me compute the ratio for each item: B21:138/594=0.232, B00:583/551=1.058, A08:868/634=1.369, A27:822/204=4.029, B94:783/557=1.406, B53:65/440=0.148, A23:262/734=0.357, C65:121/893=0.135, A89:508/246=2.065, B67:780/592=1.318, B45:461/143=3.223, C34:484/568=0.852, C12:668/472=1.415, A05:389/401=0.970, B60:808/729=1.108, C92:215/201=1.070, B69:97/85=1.141, A84:500/641=0.780, C06:30/958=0.031, B12:915/80=11.438, A16:856/157=5.451, B85:400/806=0.496, C35:444/684=0.649, B10:623/59=10.559, C15:781/31=25.194, B09:786/765=1.027, B16:3/416=0.007, C81:713/392=1.819, A58:457/428=1.068, C24:273/699=0.391, A30:739/141=5.241, A66:822/606=1.356, B75:235/612=0.384, B46:606/133=4.556, C70:105/689=0.152. Now, I should sort items by ratio descending to see which ones are most efficient for me: C15 (25.194), B12 (11.438), B10 (10.559), A16 (5.451), A30 (5.241), B46 (4.556), A27 (4.029), B45 (3.223), A89 (2.065), C81 (1.819), A08 (1.369), A66 (1.356), C12 (1.415), B94 (1.406), B67 (1.318), B60 (1.108), B00 (1.058), A58 (1.068), C92 (1.070), B69 (1.141), B09 (1.027), A05 (0.970), C34 (0.852), A84 (0.780), C35 (0.649), B85 (0.496), C24 (0.391), B75 (0.384), A23 (0.357), B21 (0.232), C70 (0.152), B53 (0.148), C65 (0.135), B16 (0.007), C06 (0.031). I should aim to include as many high{-}ratio items as possible. However, I don\textbackslash{}'t know the other participant\textbackslash{}'s preferences. Since this is a negotiation, I should start with a proposal that includes many high{-}ratio items for me, but also leave room for compromise. The total effort of all items is 594+551+634+204+557+440+734+893+246+592+143+568+472+401+729+201+85+641+958+80+157+806+684+59+31+765+416+392+428+699+141+606+612+133+689 = 20207, which is way above the limit 8170, so we have to leave out many items. I should calculate the cumulative effort for my top items. Let me try to add items in order of ratio until I hit the limit. Start with C15 (effort 31, imp 781), B12 (80, 915), B10 (59, 623), A16 (157, 856), A30 (141, 739), B46 (133, 606), A27 (204, 822), B45 (143, 461), A89 (246, 508), C81 (392, 713), A08 (634, 868), ... but effort already 31+80+59+157+141+133+204+143+246+392 = 1586, plus A08 (634) = 2220, which is still under. Then A66 (606, 822) would make 2826, C12 (472, 668) = 3298, B94 (557, 783) = 3855, B67 (592, 780) = 4447, B60 (729, 808) = 5176, B00 (551, 583) = 5727, A58 (428, 457) = 6155, C92 (201, 215) = 6356, B69 (85, 97) = 6441, B09 (765, 786) = 7206, A05 (401, 389) = 7607, C34 (568, 484) = 8175 which exceeds 8170. So I should stop before C34. Actually, 7607 + 568 = 8175 > 8170, so I can include items up to A05 with effort 7607, and then I have 8170{-}7607=563 left. I can add some smaller items. I have B69 (85) already included, but I can add more like B10 is already in, etc. Actually, let me list the items in ratio order and cumulative effort: C15:31, B12:80 (111), B10:59 (170), A16:157 (327), A30:141 (468), B46:133 (601), A27:204 (805), B45:143 (948), A89:246 (1194), C81:392 (1586), A08:634 (2220), A66:606 (2826), C12:472 (3298), B94:557 (3855), B67:592 (4447), B60:729 (5176), B00:551 (5727), A58:428 (6155), C92:201 (6356), B69:85 (6441), B09:765 (7206), A05:401 (7607). Now, remaining effort 8170{-}7607=563. I can add items with high ratio that are not included, but all high ratio are already included. Lower ratio items: maybe C35 (684) is too big, B85 (806) too big, etc. I can add B21 (594) but too big. Actually, I have B10 already which is small. Perhaps I can add C24 (699) too big. Or A84 (641) too big. Or B75 (612) too big. So maybe I can add multiple small ones: like B53 (440) but ratio low, or C70 (689) too big. Actually, I have not included C15? It is included. Perhaps I can add B16 (416) but ratio very low. Or A23 (734) too big. So it seems with 563 left, I can only add smaller items that are not in the list. I have not included C65 (893) too big, C06 (958) too big, A23 (734) too big, B53 (440) but 440<563, so I can add B53 (effort 440, imp 65) which is not good for me. Or I can add C34 (568) but 568>563. So no, I cannot add any item without exceeding. So my set with cumulative effort 7607 includes items: C15, B12, B10, A16, A30, B46, A27, B45, A89, C81, A08, A66, C12, B94, B67, B60, B00, A58, C92, B69, B09, A05. This has importance 781+915+623+856+739+606+822+461+508+713+868+822+668+783+780+808+583+457+215+97+786+389 = let\textbackslash{}'s calculate: 781+915=1696, +623=2319, +856=3175, +739=3914, +606=4520, +822=5342, +461=5803, +508=6311, +713=7024, +868=7892, +822=8714, +668=9382, +783=10165, +780=10945, +808=11753, +583=12336, +457=12793, +215=13008, +97=13105, +786=13891, +389=14280. So importance 14280. But this is without the other participant\textbackslash{}'s preferences. Since I don\textbackslash{}'t know theirs, I should propose this set as it is very good for me. However, to negotiate, I might need to compromise. I should start with this proposal.'\}} \\
\\ 
\texttt{PROPOSAL: \{'C15', 'B12', 'B10', 'A16', 'A30', 'B46', 'A27', 'B45', 'A89', 'C81', 'A08', 'A66', 'C12', 'B94', 'B67', 'B60', 'B00', 'A58', 'C92', 'B69', 'B09', 'A05'\}} \\
\texttt{ARGUMENT: \{'This proposal includes items that are highly efficient in terms of importance per unit effort, ensuring we get the most value within the effort limit. It leaves out items with very low efficiency, which helps maximize our joint benefit.'\}} \\
            }
        }
    }
    & & \\ \\

    \theutterance \stepcounter{utterance}  
    & & \multicolumn{4}{p{0.6\linewidth}}{
        \cellcolor[rgb]{0.9,0.9,0.9}{
            \makecell[{{p{\linewidth}}}]{
                \texttt{\tiny{[GM$\rangle$P2]}}
                \texttt{You are participating in a collaborative negotiation game.} \\
\\ 
\texttt{Together with another participant, you must agree on a single set of items that will be kept. Each of you has your own view of how much each item matters to you (importance). You do not know how the other participant values the items. Additionally, you are given the effort each item demands.} \\
\texttt{You may only agree on a set if the total effort of the selected items does not exceed a shared limit:} \\
\\ 
\texttt{LIMIT = 8170} \\
\\ 
\texttt{Here are the individual item effort values:} \\
\\ 
\texttt{Item effort = \{"B21": 594, "B00": 551, "A08": 634, "A27": 204, "B94": 557, "B53": 440, "A23": 734, "C65": 893, "A89": 246, "B67": 592, "B45": 143, "C34": 568, "C12": 472, "A05": 401, "B60": 729, "C92": 201, "B69": 85, "A84": 641, "C06": 958, "B12": 80, "A16": 157, "B85": 806, "C35": 684, "B10": 59, "C15": 31, "B09": 765, "B16": 416, "C81": 392, "A58": 428, "C24": 699, "A30": 141, "A66": 606, "B75": 612, "B46": 133, "C70": 689\}} \\
\\ 
\texttt{Here is your personal view on the importance of each item:} \\
\\ 
\texttt{Item importance values = \{"B21": 138, "B00": 583, "A08": 868, "A27": 822, "B94": 783, "B53": 65, "A23": 262, "C65": 121, "A89": 508, "B67": 780, "B45": 461, "C34": 484, "C12": 668, "A05": 389, "B60": 808, "C92": 215, "B69": 97, "A84": 500, "C06": 30, "B12": 915, "A16": 856, "B85": 400, "C35": 444, "B10": 623, "C15": 781, "B09": 786, "B16": 3, "C81": 713, "A58": 457, "C24": 273, "A30": 739, "A66": 822, "B75": 235, "B46": 606, "C70": 105\}} \\
\\ 
\texttt{Goal:} \\
\\ 
\texttt{Your goal is to negotiate a shared set of items that benefits you as much as possible (i.e., maximizes total importance to YOU), while staying within the LIMIT. You are not required to make a PROPOSAL in every message {-} you can simply negotiate as well. All tactics are allowed!} \\
\\ 
\texttt{Interaction Protocol:} \\
\\ 
\texttt{You may only use the following structured formats in a message:} \\
\\ 
\texttt{PROPOSAL: \{'A', 'B', 'C', …\}} \\
\texttt{Propose keeping exactly those items.} \\
\texttt{REFUSE: \{'A', 'B', 'C', …\}} \\
\texttt{Explicitly reject opponent's proposal.} \\
\texttt{ARGUMENT: \{'...'\}} \\
\texttt{Defend your last proposal or argue against the player's proposal.} \\
\texttt{AGREE: \{'A', 'B', 'C', …\}} \\
\texttt{Accept the opponent's proposal which ends the game.} \\
\texttt{STRATEGIC REASONING: \{'...'\}} \\
\texttt{	Describe your strategic reasoning or anticipation explaining your choice of action. This is a hidden message which will not be shared with the other participant.} \\
\\ 
\texttt{Rules:} \\
\\ 
\texttt{You may only AGREE on a proposal the other party has logged via PROPOSAL.} \\
\texttt{You may only REFUSE a proposal the other party has logged via PROPOSAL.} \\
\texttt{Total effort of any PROPOSAL or AGREE set must be ≤ LIMIT.} \\
\texttt{Do NOT reveal your hidden importance scores.} \\
\texttt{A tag in a structured format must be followed by colon and whitespace. The argument must be a python set containing 0 or more strings.} \\
\texttt{So, it must be of the form TAG: \{...\}} \\
\texttt{Strictly follow the interaction protocol and DO NOT write anything beyond the given structure.} \\
\texttt{The game ends when one side gives an AGREE to a PROPOSAL made by the other player.} \\
\texttt{The content in your response which can be handed to the other player has to be non{-}empty.} \\
\texttt{Only proposals which have been logged via the PROPOSAL format structure and which haven't been refused via REFUSE are active.} \\
\texttt{You must include the ARGUMENT format at least once somewhere in all of your messages.} \\
\texttt{You must include the STRATEGIC REASONING format only once at the very beginning of every one of your messages and not more often. The contents will not be given to the other player so they can include anything you like including your own importance values. Here you should reason step by step to come up with you next move.} \\
\\ 
\texttt{You will now receive the first message of the other player.} \\
\\ 
\texttt{PROPOSAL: \{'C15', 'B12', 'B10', 'A16', 'A30', 'B46', 'A27', 'B45', 'A89', 'C81', 'A08', 'A66', 'C12', 'B94', 'B67', 'B60', 'B00', 'A58', 'C92', 'B69', 'B09', 'A05'\}} \\
\texttt{ARGUMENT: \{'This proposal includes items that are highly efficient in terms of importance per unit effort, ensuring we get the most value within the effort limit. It leaves out items with very low efficiency, which helps maximize our joint benefit.'\}} \\
            }
        }
    }
    & \\ \\

    \theutterance \stepcounter{utterance}  
    & & & \multicolumn{4}{p{0.6\linewidth}}{
        \cellcolor[rgb]{0.9,0.9,0.9}{
            \makecell[{{p{\linewidth}}}]{
                \texttt{\tiny{[GM$\langle$P2]}}
                \texttt{STRATEGIC REASONING: \{'I need to analyze the opponent\textbackslash{}'s proposal from my perspective. First, I\textbackslash{}'ll calculate the total effort of their proposal to ensure it\textbackslash{}'s within the limit (8170). Then, I\textbackslash{}'ll compute the total importance to me of their proposed set. Their proposal includes: C15, B12, B10, A16, A30, B46, A27, B45, A89, C81, A08, A66, C12, B94, B67, B60, B00, A58, C92, B69, B09, A05. I have my importance values: B21:138, B00:583, A08:868, A27:822, B94:783, B53:65, A23:262, C65:121, A89:508, B67:780, B45:461, C34:484, C12:668, A05:389, B60:808, C92:215, B69:97, A84:500, C06:30, B12:915, A16:856, B85:400, C35:444, B10:623, C15:781, B09:786, B16:3, C81:713, A58:457, C24:273, A30:739, A66:822, B75:235, B46:606, C70:105. And efforts: B21:594, B00:551, A08:634, A27:204, B94:557, B53:440, A23:734, C65:893, A89:246, B67:592, B45:143, C34:568, C12:472, A05:401, B60:729, C92:201, B69:85, A84:641, C06:958, B12:80, A16:157, B85:806, C35:684, B10:59, C15:31, B09:765, B16:416, C81:392, A58:428, C24:699, A30:141, A66:606, B75:612, B46:133, C70:689.} \\
\\ 
\texttt{Calculating total effort of their proposal: } \\
\texttt{C15:31, B12:80, B10:59, A16:157, A30:141, B46:133, A27:204, B45:143, A89:246, C81:392, A08:634, A66:606, C12:472, B94:557, B67:592, B60:729, B00:551, A58:428, C92:201, B69:85, B09:765, A05:401.} \\
\texttt{Sum: 31+80=111; +59=170; +157=327; +141=468; +133=601; +204=805; +143=948; +246=1194; +392=1586; +634=2220; +606=2826; +472=3298; +557=3855; +592=4447; +729=5176; +551=5727; +428=6155; +201=6356; +85=6441; +765=7206; +401=7607.} \\
\texttt{Total effort = 7607, which is under 8170, so valid.} \\
\\ 
\texttt{Now, my importance from their set: } \\
\texttt{C15:781, B12:915, B10:623, A16:856, A30:739, B46:606, A27:822, B45:461, A89:508, C81:713, A08:868, A66:822, C12:668, B94:783, B67:780, B60:808, B00:583, A58:457, C92:215, B69:97, B09:786, A05:389.} \\
\texttt{Sum: 781+915=1696; +623=2319; +856=3175; +739=3914; +606=4520; +822=5342; +461=5803; +508=6311; +713=7024; +868=7892; +822=8714; +668=9382; +783=10165; +780=10945; +808=11753; +583=12336; +457=12793; +215=13008; +97=13105; +786=13891; +389=14280.} \\
\texttt{So total importance to me = 14280.} \\
\\ 
\texttt{This is a very high value for me. However, I should consider if I can do better by including some items they excluded and removing some with lower efficiency for me. For instance, items like A84 (importance 500, effort 641) and C34 (484, 568) might be worth adding if we remove less efficient ones. But their proposal already has many high{-}importance items for me. I should also think about what the opponent might value, as they proposed this set. They might value items I don\textbackslash{}'t care much about, so I should be cautious.} \\
\\ 
\texttt{Looking at items not in their proposal: B21 (138/594, low efficiency), B53 (65/440, very low), A23 (262/734, low), C65 (121/893, very low), C06 (30/958, terrible), B16 (3/416, terrible), C35 (444/684, medium), B85 (400/806, low), C24 (273/699, low), B75 (235/612, low), C70 (105/689, low). Most excluded items have low importance to me, except perhaps C35 (444) and maybe others, but their efficiencies are not great.} \\
\\ 
\texttt{My top items by importance per effort: } \\
\texttt{B12: 915/80 = 11.4375} \\
\texttt{C15: 781/31 = 25.1935} \\
\texttt{B10: 623/59 = 10.5593} \\
\texttt{A16: 856/157 = 5.4522} \\
\texttt{A30: 739/141 = 5.2411} \\
\texttt{B46: 606/133 = 4.5564} \\
\texttt{A27: 822/204 = 4.0294} \\
\texttt{B45: 461/143 = 3.2238} \\
\texttt{A89: 508/246 = 2.065} \\
\texttt{C81: 713/392 = 1.819} \\
\texttt{A08: 868/634 = 1.369} \\
\texttt{A66: 822/606 = 1.356} \\
\texttt{C12: 668/472 = 1.415} \\
\texttt{B94: 783/557 = 1.406} \\
\texttt{B67: 780/592 = 1.318} \\
\texttt{B60: 808/729 = 1.108} \\
\texttt{B00: 583/551 = 1.058} \\
\texttt{A58: 457/428 = 1.068} \\
\texttt{C92: 215/201 = 1.069} \\
\texttt{B69: 97/85 = 1.141} \\
\texttt{B09: 786/765 = 1.027} \\
\texttt{A05: 389/401 = 0.97} \\
\\ 
\texttt{Many are already included. Excluded items with decent efficiency: C35:444/684=0.649, which is lower than many included. A84:500/641=0.78, also lower. So their proposal is actually very good for me. I should agree to it, as it gives me 14280 importance, which is likely near maximum for me. I might not get a better deal. However, to be strategic, I should not agree immediately but perhaps propose a minor modification to see if I can get slightly more, but risk losing this good set.} \\
\\ 
\texttt{But since the opponent might have proposed this thinking it benefits them too, I should agree quickly to secure it. Arguing might lead to them changing their mind. So I will agree.'\}} \\
\\ 
\texttt{AGREE: \{'C15', 'B12', 'B10', 'A16', 'A30', 'B46', 'A27', 'B45', 'A89', 'C81', 'A08', 'A66', 'C12', 'B94', 'B67', 'B60', 'B00', 'A58', 'C92', 'B69', 'B09', 'A05'\}} \\
            }
        }
    }
     \\ \\

    \theutterance \stepcounter{utterance}  
    & & \multicolumn{4}{p{0.6\linewidth}}{
        \cellcolor[rgb]{0.9,0.9,0.9}{
            \makecell[{{p{\linewidth}}}]{
                \texttt{\tiny{[GM$\rangle$P2]}}
                \texttt{Your response did not contain an argument. You must include the structured format ARGUMENT: \{...\} somewhere in your response. Try again} \\
            }
        }
    }
    & \\ \\

    \theutterance \stepcounter{utterance}  
    & & & \multicolumn{4}{p{0.6\linewidth}}{
        \cellcolor[rgb]{0.9,0.9,0.9}{
            \makecell[{{p{\linewidth}}}]{
                \texttt{\tiny{[GM$\langle$P2]}}
                \texttt{STRATEGIC REASONING: \{'The opponent\textbackslash{}'s proposal has a total effort of 7607, which is under the limit, and gives me a very high importance value of 14280. This is likely near the maximum achievable for me, as it includes almost all high{-}efficiency items from my perspective (like B12, C15, B10, A16, A30). The excluded items (such as B53, A23, C65, C06, B16) have very low importance to me, so adding them would not significantly improve my score and would use effort that could be better spent. Since this proposal is already excellent for me, I should accept it immediately to avoid the risk of the opponent retracting it or proposing a worse set. I will include a brief argument to justify my agreement.'\}} \\
\\ 
\texttt{ARGUMENT: \{'This is a well{-}balanced proposal that efficiently uses the effort limit to include many high{-}value items, making it beneficial for both of us.'\}} \\
\texttt{AGREE: \{'C15', 'B12', 'B10', 'A16', 'A30', 'B46', 'A27', 'B45', 'A89', 'C81', 'A08', 'A66', 'C12', 'B94', 'B67', 'B60', 'B00', 'A58', 'C92', 'B69', 'B09', 'A05'\}} \\
            }
        }
    }
     \\ \\

    \theutterance \stepcounter{utterance}  
    & & & \multicolumn{2}{p{0.3\linewidth}}{
        \cellcolor[rgb]{0.9,0.9,0.9}{
            \makecell[{{p{\linewidth}}}]{
                \texttt{\tiny{[GM$|$GM]}}
                \texttt{game successful} \\
            }
        }
    }
    & & \\ \\

    \theutterance \stepcounter{utterance}  
    & & & \multicolumn{2}{p{0.3\linewidth}}{
        \cellcolor[rgb]{0.9,0.9,0.9}{
            \makecell[{{p{\linewidth}}}]{
                \texttt{\tiny{[GM$|$GM]}}
                \texttt{end game} \\
            }
        }
    }
    & & \\ \\

\end{supertabular}
}

\end{document}
